\section{Scope, Assumptions, and Falsifiability}
\label{sec:scope-falsifiability}

\subsection{Scope of the Preprint}

This preprint has a deliberately narrow and explicitly delimited scope.
It is a diagnostic analysis of \emph{KPI-based observation} for enterprise
continua modeled under the Ontology of Continua (OC) and the frozen executable
theory specification \texttt{ETS.K\_EA.v1.1}.

We do not propose a new enterprise model, a new KPI methodology, nor a new
estimation, optimization, governance, or decision framework.
Instead, we make explicit structural consequences that already follow from the
OC/ETS setup itself:

\begin{itemize}[leftmargin=2em]
  \item the enterprise is represented as an OC/ETS continuum with state space
        \(\Omega(K_{EA})\) and an admissible observable domain
        \(\Omega_{\mathrm{obs}}\subseteq\Omega(K_{EA})\);
  \item KPI observation is modeled as a finite readout map
        \(\pi_K:\Omega_{\mathrm{obs}}\to\Sigma_{\mathrm{obs}}\subseteq\mathbb{R}^m\);
  \item identifiability is defined strictly as injectivity of \(\pi_K\) on a
        specified domain of interest.
\end{itemize}

The central conclusion is not that KPIs are ineffective, misleading, or
inappropriate, but that identifiability cannot be inferred from OC/ETS alone,
and that deterministic post-processing cannot recover information already lost
at the level of KPI readout.

\subsection{Explicit Assumptions}

All results in this preprint rely exclusively on the following explicit and
exhaustive assumptions:

\begin{enumerate}[leftmargin=2em]
  \item \textbf{OC/ETS state space.}
        The enterprise continuum \(K_{EA}\) is described by OC primitives
        (axes, flows, thresholds, cycles, and a continuum measure), with state
        space \(\Omega(K_{EA})\).
        Observation is restricted to the admissible domain
        \(\Omega_{\mathrm{obs}}\), defined solely by ETS-level observability
        structures.
  \item \textbf{Finite KPI observation.}
        A KPI family \(K=(\kappa_1,\ldots,\kappa_m)\) is finite and admissible
        under \texttt{ETS.K\_EA.v1.1}, inducing a readout map
        \(\pi_K:\Omega_{\mathrm{obs}}\to\Sigma_{\mathrm{obs}}\subseteq\mathbb{R}^m\).
  \item \textbf{No hidden enrichment.}
        Any analytic, statistical, or computational procedure that operates only
        on KPI values is modeled as a deterministic post-processing map composed
        with \(\pi_K\).
\end{enumerate}

No assumptions of smoothness, continuity, dimensionality, genericity,
regularity, noise structure, or probabilistic modeling are required for the
logical core of the results in Section~\ref{sec:impossibility}.

\subsection{What Is Explicitly Not Claimed}

This preprint explicitly does \emph{not} claim:

\begin{itemize}[leftmargin=2em]
  \item that KPI observation cannot be identifying on restricted domains
        \(D\subseteq\Omega_{\mathrm{obs}}\);
  \item that identifiability cannot be achieved by introducing additional
        observables, sensors, experimental interventions, or structural
        constraints beyond OC/ETS;
  \item that any concrete KPI system used in practice is necessarily
        non-identifying or analytically deficient.
\end{itemize}

The claim is strictly limited to the following statement:
OC and \texttt{ETS.K\_EA.v1.1} alone do not entail injectivity of an arbitrary
finite KPI readout map, and deterministic post-processing cannot recover
information lost within KPI fibers.

\subsection{Falsifiability and Empirical Interface}

Because the results of this preprint are diagnostic and largely negative,
falsifiability must be interpreted at the appropriate formal and empirical
interfaces.

\begin{enumerate}[leftmargin=2em]
  \item \textbf{Formal falsification (model level).}
        If OC or \texttt{ETS.K\_EA.v1.1} were extended by an explicit axiom
        enforcing injectivity of \(\pi_K\) on \(\Omega_{\mathrm{obs}}\) for a
        specified class of KPI families, then
        Theorem~\ref{thm:no-complete} would no longer hold within that amended
        formal system.
        This would constitute a formal falsification of the scope of the present
        results.
  \item \textbf{Empirical falsification (measurement level).}
        For a fixed enterprise, a fixed KPI family \(K\), and a fixed readout
        protocol, non-identifiability can be empirically witnessed by exhibiting
        two distinct operationally realizable states
        \(\omega_1\neq\omega_2\in\Omega_{\mathrm{obs}}\) such that
        \(\pi_K(\omega_1)=\pi_K(\omega_2)\).
        Existence of such a pair directly demonstrates non-identifiability of
        \(K\) on the tested domain.
        Conversely, failure to find such pairs does not establish global
        injectivity; it only supports identifiability for the explored domain and
        protocol.
  \item \textbf{Protocol dependence.}
        Because KPI values depend on sampling, aggregation windows, and cadence,
        any empirical claim must fix the readout protocol explicitly.
        Once fixed, the protocol is absorbed into the definition of \(\pi_K\),
        and the formal results apply unchanged.
\end{enumerate}

Accordingly, the falsifiable content of this preprint can be summarized as
follows:
\emph{for a fixed KPI readout protocol, if two distinct admissible enterprise
states yield identical KPI values, then KPI-based observation is
non-identifying on that tested domain}.
This statement is a direct consequence of the definitions in
Section~\ref{sec:formal-setting}.
